%%%%%%%%%%%%%%%%%%%%%%%%%%%%%%%%%%%%%%%%%%%%%%%%%%%%%%%%%%%%
% 0.6 FUNCTIONS
% The Composition of Advanced Mathematics
%%%%%%%%%%%%%%%%%%%%%%%%%%%%%%%%%%%%%%%%%%%%%%%%%%%%%%%%%%%%

\section{Functions}
\label{sec:functions}

\prerequisite{
Familiarity with sets from \cref{sec:elementary-set-theory} and relations
from \cref{sec:relations}.
}

\learningobjectives{
By the end of this section, the reader should be able to:
\begin{itemize}
    \item define a function as a special type of relation;
    \item distinguish domain, codomain, and range;
    \item interpret function notation and mapping notation;
    \item evaluate functions and determine equality of functions;
    \item recognize identity, constant, inclusion, and piecewise-defined functions;
    \item determine whether a function is injective, surjective, or bijective;
    \item compose functions and interpret the order of composition;
    \item understand inverse functions and the connection between invertibility
          and bijectivity;
    \item work with images and preimages of sets;
    \item interpret sequences and multivariable functions as functions;
    \item recognize common errors involving functions.
\end{itemize}
}

%%%%%%%%%%%%%%%%%%%%%%%%%%%%%%%%%%%%%%%%%%%%%%%%%%%%%%%%%%%%
\subsection{Why Functions Matter}
\label{subsec:why-functions-matter}
%%%%%%%%%%%%%%%%%%%%%%%%%%%%%%%%%%%%%%%%%%%%%%%%%%%%%%%%%%%%

Functions are among the most central objects in mathematics.

They describe how one quantity, object, or structure depends on another.

Examples include:

\begin{itemize}
    \item the area of a circle as a function of its radius;
    \item position as a function of time;
    \item temperature as a function of location;
    \item probability as a function of an event;
    \item a matrix transformation acting on vectors;
    \item a derivative assigning a new function to a differentiable function;
    \item a statistical model mapping data to predictions.
\end{itemize}

Nearly every branch of mathematics uses functions in some form.

%%%%%%%%%%%%%%%%%%%%%%%%%%%%%%%%%%%%%%%%%%%%%%%%%%%%%%%%%%%%
\subsection{Functions as Special Relations}
\label{subsec:functions-special-relations}
%%%%%%%%%%%%%%%%%%%%%%%%%%%%%%%%%%%%%%%%%%%%%%%%%%%%%%%%%%%%

Recall that a relation from \(A\) to \(B\) is a subset of

\[ A\times B. \]

A function is a relation with an additional requirement.

\begin{definitionbox}{Function}
A \emph{function} \(f\) from a set \(A\) to a set \(B\) is a relation that
associates each element of \(A\) with exactly one element of \(B\).
\end{definitionbox}

Thus every input must have one and only one output.

%%%%%%%%%%%%%%%%%%%%%%%%%%%%%%%%%%%%%%%%%%%%%%%%%%%%%%%%%%%%
\subsection{Function Mapping Notation}
\label{subsec:function-mapping-notation}
%%%%%%%%%%%%%%%%%%%%%%%%%%%%%%%%%%%%%%%%%%%%%%%%%%%%%%%%%%%%

A function from \(A\) to \(B\) is commonly written

\[ f:A\to B. \]

The arrow

\[ \to \]

is read

\begin{center}
``maps to''
\end{center}

in this context.

Thus,

\[ f:A\to B \]

is read

\begin{center}
``\(f\) maps \(A\) to \(B\)''
\end{center}

or

\begin{center}
``\(f\) is a function from \(A\) to \(B\).''
\end{center}

The same arrow has different readings in other contexts. For example,
\(x\to a\) in calculus is read ``\(x\) approaches \(a\).'' Context determines
the intended meaning.

%%%%%%%%%%%%%%%%%%%%%%%%%%%%%%%%%%%%%%%%%%%%%%%%%%%%%%%%%%%%
\subsection{The Mapsto Symbol}
\label{subsec:mapsto-symbol}
%%%%%%%%%%%%%%%%%%%%%%%%%%%%%%%%%%%%%%%%%%%%%%%%%%%%%%%%%%%%

The symbol

\[ \mapsto \]

is called the \emph{mapsto symbol} and is read

\begin{center}
``maps to.''
\end{center}

For example,

\[ x\mapsto x^2 \]

means that the input \(x\) is assigned the output \(x^2\).

A complete description may be written

\[ f:\R\to\R,\qquad x\mapsto x^2. \]

%%%%%%%%%%%%%%%%%%%%%%%%%%%%%%%%%%%%%%%%%%%%%%%%%%%%%%%%%%%%
\subsection{Function Notation}
\label{subsec:function-notation}
%%%%%%%%%%%%%%%%%%%%%%%%%%%%%%%%%%%%%%%%%%%%%%%%%%%%%%%%%%%%

If \(f\) is a function and \(x\) is an input, the corresponding output is
written

\[ f(x). \]

This is read

\begin{center}
``\(f\) of \(x\).''
\end{center}

The parentheses do not mean multiplication.

Thus,

\[ f(x) \]

does not mean

\[ f\cdot x. \]

%%%%%%%%%%%%%%%%%%%%%%%%%%%%%%%%%%%%%%%%%%%%%%%%%%%%%%%%%%%%
\subsection{Domain}
\label{subsec:function-domain}
%%%%%%%%%%%%%%%%%%%%%%%%%%%%%%%%%%%%%%%%%%%%%%%%%%%%%%%%%%%%

\begin{definitionbox}{Domain}
For a function
\[ f:A\to B, \]
the set \(A\) is called the \emph{domain} of \(f\).
\end{definitionbox}

The domain is the set of allowed inputs.

It may be written

\[ \dom(f)=A. \]

%%%%%%%%%%%%%%%%%%%%%%%%%%%%%%%%%%%%%%%%%%%%%%%%%%%%%%%%%%%%
\subsection{Codomain}
\label{subsec:function-codomain}
%%%%%%%%%%%%%%%%%%%%%%%%%%%%%%%%%%%%%%%%%%%%%%%%%%%%%%%%%%%%

\begin{definitionbox}{Codomain}
For a function
\[ f:A\to B, \]
the set \(B\) is called the \emph{codomain} of \(f\).
\end{definitionbox}

The codomain is the set in which the outputs are required to lie.

It may contain elements that are never actually produced by the function.

%%%%%%%%%%%%%%%%%%%%%%%%%%%%%%%%%%%%%%%%%%%%%%%%%%%%%%%%%%%%
\subsection{Range or Image}
\label{subsec:function-range}
%%%%%%%%%%%%%%%%%%%%%%%%%%%%%%%%%%%%%%%%%%%%%%%%%%%%%%%%%%%%

\begin{definitionbox}{Range}
The \emph{range}, or \emph{image}, of a function \(f:A\to B\) is the set of
actual outputs:
\[ \Image(f)=\{f(x):x\in A\}. \]
\end{definitionbox}

The notation

\[ \Image(f) \]

is read

\begin{center}
``the image of \(f\).''
\end{center}

The range always satisfies

\[ \Image(f)\subseteq B. \]

%%%%%%%%%%%%%%%%%%%%%%%%%%%%%%%%%%%%%%%%%%%%%%%%%%%%%%%%%%%%
\subsection{Domain, Codomain, and Range}
\label{subsec:domain-codomain-range}
%%%%%%%%%%%%%%%%%%%%%%%%%%%%%%%%%%%%%%%%%%%%%%%%%%%%%%%%%%%%

These three concepts should not be confused.

For

\[ f:A\to B, \]

we have:

\begin{itemize}
    \item \(A\): domain;
    \item \(B\): codomain;
    \item \(\Image(f)\): actual outputs.
\end{itemize}

In general,

\[ \Image(f)\subseteq B, \]

and equality need not hold.

%%%%%%%%%%%%%%%%%%%%%%%%%%%%%%%%%%%%%%%%%%%%%%%%%%%%%%%%%%%%
\subsection{Example of Domain, Codomain, and Range}
\label{subsec:domain-codomain-range-example}
%%%%%%%%%%%%%%%%%%%%%%%%%%%%%%%%%%%%%%%%%%%%%%%%%%%%%%%%%%%%

Define

\[ f:\R\to\R,\qquad f(x)=x^2. \]

The domain is

\[ \R. \]

The codomain is also

\[ \R. \]

But the range is

\[ [0,\infty). \]

Therefore,

\[ \Image(f)\neq\R. \]

%%%%%%%%%%%%%%%%%%%%%%%%%%%%%%%%%%%%%%%%%%%%%%%%%%%%%%%%%%%%
\subsection{Evaluating Functions}
\label{subsec:evaluating-functions}
%%%%%%%%%%%%%%%%%%%%%%%%%%%%%%%%%%%%%%%%%%%%%%%%%%%%%%%%%%%%

To evaluate a function, replace its input variable by the specified value.

Suppose

\[ f(x)=3x^2-2x+1. \]

Then

\[ f(2)=3(2)^2-2(2)+1=12-4+1=9. \]

%%%%%%%%%%%%%%%%%%%%%%%%%%%%%%%%%%%%%%%%%%%%%%%%%%%%%%%%%%%%
\subsection{Worked Example: Function Evaluation}
\label{subsec:worked-function-evaluation}
%%%%%%%%%%%%%%%%%%%%%%%%%%%%%%%%%%%%%%%%%%%%%%%%%%%%%%%%%%%%

\begin{workedexamplebox}{Evaluate a Function}

Let

\[ f(x)=2x^2-5x+3. \]

Find \(f(-2)\).

\begin{tabularx}{\textwidth}{@{}>{\raggedright\arraybackslash}p{0.42\textwidth}X@{}}
Substitute \(-2\) for \(x\). & \(f(-2)=2(-2)^2-5(-2)+3\)\\
Evaluate the square. & \(2(4)+10+3\)\\
Simplify. & \(8+10+3=21\)
\end{tabularx}

\begin{answerbox}
\[ f(-2)=21 \]
\end{answerbox}
\end{workedexamplebox}

%%%%%%%%%%%%%%%%%%%%%%%%%%%%%%%%%%%%%%%%%%%%%%%%%%%%%%%%%%%%
\subsection{Functions Need Not Be Given by Formulas}
\label{subsec:functions-not-only-formulas}
%%%%%%%%%%%%%%%%%%%%%%%%%%%%%%%%%%%%%%%%%%%%%%%%%%%%%%%%%%%%

A function may be described by:

\begin{itemize}
    \item a formula;
    \item a table;
    \item a graph;
    \item an arrow diagram;
    \item an algorithm;
    \item a verbal rule;
    \item a set of ordered pairs.
\end{itemize}

The essential requirement is not the presence of a formula.

The essential requirement is that each input have exactly one output.

%%%%%%%%%%%%%%%%%%%%%%%%%%%%%%%%%%%%%%%%%%%%%%%%%%%%%%%%%%%%
\subsection{Functions as Sets of Ordered Pairs}
\label{subsec:function-ordered-pairs}
%%%%%%%%%%%%%%%%%%%%%%%%%%%%%%%%%%%%%%%%%%%%%%%%%%%%%%%%%%%%

A function may be represented as a set of ordered pairs.

For example,

\[ f=\{(1,a),(2,b),(3,b)\} \]

defines a function from

\[ \{1,2,3\} \]

to a codomain containing \(a\) and \(b\).

The first coordinates are inputs and the second coordinates are outputs.

%%%%%%%%%%%%%%%%%%%%%%%%%%%%%%%%%%%%%%%%%%%%%%%%%%%%%%%%%%%%
\subsection{When a Relation Is Not a Function}
\label{subsec:not-a-function}
%%%%%%%%%%%%%%%%%%%%%%%%%%%%%%%%%%%%%%%%%%%%%%%%%%%%%%%%%%%%

Consider

\[ R=\{(1,a),(1,b),(2,c)\}. \]

The input \(1\) is associated with two distinct outputs:

\[ a\qquad\text{and}\qquad b. \]

Therefore \(R\) is not a function.

By contrast,

\[ \{(1,a),(2,a),(3,a)\} \]

is a function because several inputs may share the same output.

%%%%%%%%%%%%%%%%%%%%%%%%%%%%%%%%%%%%%%%%%%%%%%%%%%%%%%%%%%%%
\subsection{Exactly One Output}
\label{subsec:exactly-one-output}
%%%%%%%%%%%%%%%%%%%%%%%%%%%%%%%%%%%%%%%%%%%%%%%%%%%%%%%%%%%%

The phrase ``exactly one'' combines two requirements:

\begin{enumerate}
    \item every input has at least one output;
    \item no input has more than one output.
\end{enumerate}

Both conditions are necessary.

%%%%%%%%%%%%%%%%%%%%%%%%%%%%%%%%%%%%%%%%%%%%%%%%%%%%%%%%%%%%
\subsection{Equality of Functions}
\label{subsec:function-equality}
%%%%%%%%%%%%%%%%%%%%%%%%%%%%%%%%%%%%%%%%%%%%%%%%%%%%%%%%%%%%

Two functions are equal only when they have the same domain, the same
codomain, and assign the same output to every input.

Thus,

\[ f=g \]

means that the functions agree as complete mappings, not merely that their
formulas look similar.

For example,

\[ f:\R\to\R,\qquad f(x)=x^2 \]

and

\[ g:[0,\infty)\to\R,\qquad g(x)=x^2 \]

are not the same function because their domains differ.

%%%%%%%%%%%%%%%%%%%%%%%%%%%%%%%%%%%%%%%%%%%%%%%%%%%%%%%%%%%%
\subsection{Identity Functions}
\label{subsec:identity-functions}
%%%%%%%%%%%%%%%%%%%%%%%%%%%%%%%%%%%%%%%%%%%%%%%%%%%%%%%%%%%%

\begin{definitionbox}{Identity Function}
For a set \(A\), the \emph{identity function} on \(A\) is
\[ \operatorname{id}_A:A\to A,\qquad \operatorname{id}_A(x)=x. \]
\end{definitionbox}

The notation

\[ \operatorname{id}_A \]

is read

\begin{center}
``the identity on \(A\).''
\end{center}

The identity function leaves every element unchanged.

%%%%%%%%%%%%%%%%%%%%%%%%%%%%%%%%%%%%%%%%%%%%%%%%%%%%%%%%%%%%
\subsection{Constant Functions}
\label{subsec:constant-functions}
%%%%%%%%%%%%%%%%%%%%%%%%%%%%%%%%%%%%%%%%%%%%%%%%%%%%%%%%%%%%

\begin{definitionbox}{Constant Function}
A function \(f:A\to B\) is \emph{constant} if there exists a fixed
\(c\in B\) such that
\[ f(x)=c \]
for every \(x\in A\).
\end{definitionbox}

For example,

\[ f(x)=5 \]

assigns the same output to every input.

%%%%%%%%%%%%%%%%%%%%%%%%%%%%%%%%%%%%%%%%%%%%%%%%%%%%%%%%%%%%
\subsection{Inclusion Functions}
\label{subsec:inclusion-functions}
%%%%%%%%%%%%%%%%%%%%%%%%%%%%%%%%%%%%%%%%%%%%%%%%%%%%%%%%%%%%

If

\[ A\subseteq B, \]

the inclusion function is

\[ \iota:A\to B,\qquad \iota(a)=a. \]

The Greek letter

\[ \iota \]

is called \emph{iota}, pronounced ``eye-OH-tuh.''

The inclusion function sends each element of the smaller set to the same
element viewed inside the larger set.

%%%%%%%%%%%%%%%%%%%%%%%%%%%%%%%%%%%%%%%%%%%%%%%%%%%%%%%%%%%%
\subsection{Piecewise-Defined Functions}
\label{subsec:piecewise-functions}
%%%%%%%%%%%%%%%%%%%%%%%%%%%%%%%%%%%%%%%%%%%%%%%%%%%%%%%%%%%%

A function may use different formulas on different parts of its domain.

For example,

\[ f(x)=
\begin{cases}
x^2,&x\geq0,\\
-x,&x<0.
\end{cases} \]

This is called a \emph{piecewise-defined function}.

The cases environment is read according to the conditions:

\begin{center}
``\(f(x)=x^2\) if \(x\geq0\), and \(f(x)=-x\) if \(x<0\).''
\end{center}

%%%%%%%%%%%%%%%%%%%%%%%%%%%%%%%%%%%%%%%%%%%%%%%%%%%%%%%%%%%%
\subsection{Evaluating Piecewise Functions}
\label{subsec:evaluating-piecewise}
%%%%%%%%%%%%%%%%%%%%%%%%%%%%%%%%%%%%%%%%%%%%%%%%%%%%%%%%%%%%

For the function

\[ f(x)=
\begin{cases}
x^2,&x\geq0,\\
-x,&x<0,
\end{cases} \]

we have

\[ f(3)=3^2=9 \]

because \(3\geq0\).

But

\[ f(-3)=-(-3)=3 \]

because \(-3<0\).

%%%%%%%%%%%%%%%%%%%%%%%%%%%%%%%%%%%%%%%%%%%%%%%%%%%%%%%%%%%%
\subsection{Restricted Domains}
\label{subsec:restricted-domains}
%%%%%%%%%%%%%%%%%%%%%%%%%%%%%%%%%%%%%%%%%%%%%%%%%%%%%%%%%%%%

Sometimes a formula becomes useful as a function only after its domain is
restricted.

For example,

\[ f(x)=\frac1x \]

cannot have \(0\) in its domain.

A natural real domain is therefore

\[ \R\setminus\{0\}. \]

Similarly,

\[ f(x)=\sqrt x \]

has real-valued outputs only when

\[ x\geq0. \]

Thus a natural real domain is

\[ [0,\infty). \]

%%%%%%%%%%%%%%%%%%%%%%%%%%%%%%%%%%%%%%%%%%%%%%%%%%%%%%%%%%%%
\subsection{Natural Domains}
\label{subsec:natural-domains}
%%%%%%%%%%%%%%%%%%%%%%%%%%%%%%%%%%%%%%%%%%%%%%%%%%%%%%%%%%%%

When a formula is given without an explicit domain, the \emph{natural domain}
is often taken to be the largest standard set of inputs for which the formula
is meaningful in the chosen number system.

For real-valued functions, examples include:

\[ f(x)=\frac1{x-2} \]

with natural domain

\[ \R\setminus\{2\}, \]

and

\[ g(x)=\sqrt{x+4} \]

with natural domain

\[ [-4,\infty). \]

%%%%%%%%%%%%%%%%%%%%%%%%%%%%%%%%%%%%%%%%%%%%%%%%%%%%%%%%%%%%
\subsection{Injective Functions}
\label{subsec:injective-functions}
%%%%%%%%%%%%%%%%%%%%%%%%%%%%%%%%%%%%%%%%%%%%%%%%%%%%%%%%%%%%

\begin{definitionbox}{Injective Function}
A function \(f:A\to B\) is \emph{injective}, or \emph{one-to-one}, if
\[ f(x_1)=f(x_2)\Longrightarrow x_1=x_2. \]
\end{definitionbox}

This means distinct inputs cannot produce the same output.

An equivalent statement is

\[ x_1\neq x_2\Longrightarrow f(x_1)\neq f(x_2). \]

%%%%%%%%%%%%%%%%%%%%%%%%%%%%%%%%%%%%%%%%%%%%%%%%%%%%%%%%%%%%
\subsection{Example of an Injective Function}
\label{subsec:injective-example}
%%%%%%%%%%%%%%%%%%%%%%%%%%%%%%%%%%%%%%%%%%%%%%%%%%%%%%%%%%%%

Consider

\[ f:\R\to\R,\qquad f(x)=2x+3. \]

Suppose

\[ f(x_1)=f(x_2). \]

Then

\[ 2x_1+3=2x_2+3. \]

Subtracting \(3\) gives

\[ 2x_1=2x_2. \]

Dividing by \(2\) gives

\[ x_1=x_2. \]

Therefore \(f\) is injective.

%%%%%%%%%%%%%%%%%%%%%%%%%%%%%%%%%%%%%%%%%%%%%%%%%%%%%%%%%%%%
\subsection{A Function That Is Not Injective}
\label{subsec:not-injective-example}
%%%%%%%%%%%%%%%%%%%%%%%%%%%%%%%%%%%%%%%%%%%%%%%%%%%%%%%%%%%%

Consider

\[ f:\R\to\R,\qquad f(x)=x^2. \]

Since

\[ f(2)=4 \]

and

\[ f(-2)=4, \]

distinct inputs produce the same output.

Therefore \(f\) is not injective.

%%%%%%%%%%%%%%%%%%%%%%%%%%%%%%%%%%%%%%%%%%%%%%%%%%%%%%%%%%%%
\subsection{Surjective Functions}
\label{subsec:surjective-functions}
%%%%%%%%%%%%%%%%%%%%%%%%%%%%%%%%%%%%%%%%%%%%%%%%%%%%%%%%%%%%

\begin{definitionbox}{Surjective Function}
A function \(f:A\to B\) is \emph{surjective}, or \emph{onto}, if every
element of the codomain is the output of at least one input.
\end{definitionbox}

Symbolically,

\[ \forall y\in B,\;\exists x\in A\text{ such that }f(x)=y. \]

Thus,

\[ \Image(f)=B. \]

%%%%%%%%%%%%%%%%%%%%%%%%%%%%%%%%%%%%%%%%%%%%%%%%%%%%%%%%%%%%
\subsection{Example of a Surjective Function}
\label{subsec:surjective-example}
%%%%%%%%%%%%%%%%%%%%%%%%%%%%%%%%%%%%%%%%%%%%%%%%%%%%%%%%%%%%

Consider

\[ f:\R\to\R,\qquad f(x)=x^3. \]

Given any

\[ y\in\R, \]

choose

\[ x=\sqrt[3]{y}. \]

Then

\[ f(x)=\left(\sqrt[3]{y}\right)^3=y. \]

Therefore every real number has a preimage, and \(f\) is surjective.

%%%%%%%%%%%%%%%%%%%%%%%%%%%%%%%%%%%%%%%%%%%%%%%%%%%%%%%%%%%%
\subsection{Codomain Matters for Surjectivity}
\label{subsec:codomain-surjectivity}
%%%%%%%%%%%%%%%%%%%%%%%%%%%%%%%%%%%%%%%%%%%%%%%%%%%%%%%%%%%%

Consider

\[ f:\R\to\R,\qquad f(x)=x^2. \]

This function is not surjective because negative real numbers are never
outputs.

However,

\[ g:\R\to[0,\infty),\qquad g(x)=x^2 \]

is surjective.

Thus surjectivity depends on the specified codomain.

%%%%%%%%%%%%%%%%%%%%%%%%%%%%%%%%%%%%%%%%%%%%%%%%%%%%%%%%%%%%
\subsection{Bijective Functions}
\label{subsec:bijective-functions}
%%%%%%%%%%%%%%%%%%%%%%%%%%%%%%%%%%%%%%%%%%%%%%%%%%%%%%%%%%%%

\begin{definitionbox}{Bijective Function}
A function is \emph{bijective}, or a \emph{bijection}, if it is both
injective and surjective.
\end{definitionbox}

A bijection pairs the elements of the domain and codomain without omissions
or repetitions.

Each output has exactly one corresponding input.

%%%%%%%%%%%%%%%%%%%%%%%%%%%%%%%%%%%%%%%%%%%%%%%%%%%%%%%%%%%%
\subsection{Worked Example: Testing Bijectivity}
\label{subsec:worked-bijectivity}
%%%%%%%%%%%%%%%%%%%%%%%%%%%%%%%%%%%%%%%%%%%%%%%%%%%%%%%%%%%%

\begin{workedexamplebox}{Determine Whether a Function Is Bijective}

Consider

\[ f:\R\to\R,\qquad f(x)=3x-5. \]

\begin{tabularx}{\textwidth}{@{}>{\raggedright\arraybackslash}p{0.42\textwidth}X@{}}
Test injectivity. & \(3x_1-5=3x_2-5\Rightarrow x_1=x_2\)\\
Take arbitrary \(y\in\R\). & \(y=3x-5\)\\
Solve for \(x\). & \(x=\frac{y+5}{3}\in\R\)\\
Therefore every \(y\) has a preimage. & \(f\) is surjective
\end{tabularx}

\begin{answerbox}
\(f\) is bijective.
\end{answerbox}
\end{workedexamplebox}

%%%%%%%%%%%%%%%%%%%%%%%%%%%%%%%%%%%%%%%%%%%%%%%%%%%%%%%%%%%%
\subsection{Composition of Functions}
\label{subsec:function-composition}
%%%%%%%%%%%%%%%%%%%%%%%%%%%%%%%%%%%%%%%%%%%%%%%%%%%%%%%%%%%%

Suppose

\[ f:A\to B \]

and

\[ g:B\to C. \]

Their composition is

\[ g\circ f:A\to C. \]

The expression

\[ g\circ f \]

is read

\begin{center}
``\(g\) composed with \(f\).''
\end{center}

The definition is

\[ (g\circ f)(x)=g(f(x)). \]

The function \(f\) acts first, then \(g\).

%%%%%%%%%%%%%%%%%%%%%%%%%%%%%%%%%%%%%%%%%%%%%%%%%%%%%%%%%%%%
\subsection{Order Matters in Composition}
\label{subsec:composition-order}
%%%%%%%%%%%%%%%%%%%%%%%%%%%%%%%%%%%%%%%%%%%%%%%%%%%%%%%%%%%%

In general,

\[ g\circ f\neq f\circ g. \]

For example, let

\[ f(x)=x+1,\qquad g(x)=x^2. \]

Then

\[ (g\circ f)(x)=(x+1)^2, \]

while

\[ (f\circ g)(x)=x^2+1. \]

These are different functions.

%%%%%%%%%%%%%%%%%%%%%%%%%%%%%%%%%%%%%%%%%%%%%%%%%%%%%%%%%%%%
\subsection{Associativity of Composition}
\label{subsec:composition-associativity}
%%%%%%%%%%%%%%%%%%%%%%%%%%%%%%%%%%%%%%%%%%%%%%%%%%%%%%%%%%%%

Function composition is associative.

For compatible functions,

\[ h\circ(g\circ f)=(h\circ g)\circ f. \]

Indeed, both sides send \(x\) to

\[ h(g(f(x))). \]

Thus parentheses may be omitted in long compositions when the order is clear.

%%%%%%%%%%%%%%%%%%%%%%%%%%%%%%%%%%%%%%%%%%%%%%%%%%%%%%%%%%%%
\subsection{Identity and Composition}
\label{subsec:identity-composition}
%%%%%%%%%%%%%%%%%%%%%%%%%%%%%%%%%%%%%%%%%%%%%%%%%%%%%%%%%%%%

If

\[ f:A\to B, \]

then

\[ \operatorname{id}_B\circ f=f \]

and

\[ f\circ\operatorname{id}_A=f. \]

Thus identity functions behave as identity elements for function composition.

%%%%%%%%%%%%%%%%%%%%%%%%%%%%%%%%%%%%%%%%%%%%%%%%%%%%%%%%%%%%
\subsection{Worked Example: Function Composition}
\label{subsec:worked-function-composition}
%%%%%%%%%%%%%%%%%%%%%%%%%%%%%%%%%%%%%%%%%%%%%%%%%%%%%%%%%%%%

\begin{workedexamplebox}{Compose Two Functions}

Let

\[ f(x)=2x+1,\qquad g(x)=x^2. \]

Find \(g\circ f\).

\begin{tabularx}{\textwidth}{@{}>{\raggedright\arraybackslash}p{0.42\textwidth}X@{}}
Use the composition definition. & \((g\circ f)(x)=g(f(x))\)\\
Substitute \(f(x)=2x+1\) into \(g\). & \(g(2x+1)=(2x+1)^2\)\\
Expand if desired. & \(4x^2+4x+1\)
\end{tabularx}

\begin{answerbox}
\[ (g\circ f)(x)=(2x+1)^2=4x^2+4x+1 \]
\end{answerbox}
\end{workedexamplebox}

%%%%%%%%%%%%%%%%%%%%%%%%%%%%%%%%%%%%%%%%%%%%%%%%%%%%%%%%%%%%
\subsection{Inverse Functions}
\label{subsec:inverse-functions}
%%%%%%%%%%%%%%%%%%%%%%%%%%%%%%%%%%%%%%%%%%%%%%%%%%%%%%%%%%%%

The notation

\[ f^{-1} \]

is read

\begin{center}
``\(f\) inverse.''
\end{center}

\begin{definitionbox}{Inverse Function}
If \(f:A\to B\), an \emph{inverse function} is a function
\[ f^{-1}:B\to A \]
such that
\[ f^{-1}\circ f=\operatorname{id}_A \]
and
\[ f\circ f^{-1}=\operatorname{id}_B. \]
\end{definitionbox}

The inverse reverses the action of \(f\).

%%%%%%%%%%%%%%%%%%%%%%%%%%%%%%%%%%%%%%%%%%%%%%%%%%%%%%%%%%%%
\subsection{Inverse Functions Are Not Reciprocals}
\label{subsec:inverse-not-reciprocal}
%%%%%%%%%%%%%%%%%%%%%%%%%%%%%%%%%%%%%%%%%%%%%%%%%%%%%%%%%%%%

The notation

\[ f^{-1}(x) \]

does not mean

\[ \frac1{f(x)}. \]

These are entirely different operations.

For example, if

\[ f(x)=2x+3, \]

then

\[ f^{-1}(x)=\frac{x-3}{2}. \]

But

\[ \frac1{f(x)}=\frac1{2x+3}. \]

%%%%%%%%%%%%%%%%%%%%%%%%%%%%%%%%%%%%%%%%%%%%%%%%%%%%%%%%%%%%
\subsection{Finding an Inverse Algebraically}
\label{subsec:finding-inverse-algebraically}
%%%%%%%%%%%%%%%%%%%%%%%%%%%%%%%%%%%%%%%%%%%%%%%%%%%%%%%%%%%%

To find the inverse of a bijective real-valued function given by a formula, a
common procedure is:

\begin{enumerate}
    \item write \(y=f(x)\);
    \item solve the equation for \(x\) in terms of \(y\);
    \item interchange the variable names;
    \item verify the compositions if needed.
\end{enumerate}

%%%%%%%%%%%%%%%%%%%%%%%%%%%%%%%%%%%%%%%%%%%%%%%%%%%%%%%%%%%%
\subsection{Worked Example: Finding an Inverse}
\label{subsec:worked-finding-inverse}
%%%%%%%%%%%%%%%%%%%%%%%%%%%%%%%%%%%%%%%%%%%%%%%%%%%%%%%%%%%%

\begin{workedexamplebox}{Find an Inverse Function}

Let

\[ f:\R\to\R,\qquad f(x)=4x-7. \]

\begin{tabularx}{\textwidth}{@{}>{\raggedright\arraybackslash}p{0.42\textwidth}X@{}}
Write the output as \(y\). & \(y=4x-7\)\\
Add \(7\). & \(y+7=4x\)\\
Divide by \(4\). & \(x=\frac{y+7}{4}\)\\
Rename the input variable. & \(f^{-1}(x)=\frac{x+7}{4}\)
\end{tabularx}

\begin{answerbox}
\[ f^{-1}(x)=\frac{x+7}{4} \]
\end{answerbox}
\end{workedexamplebox}

%%%%%%%%%%%%%%%%%%%%%%%%%%%%%%%%%%%%%%%%%%%%%%%%%%%%%%%%%%%%
\subsection{Invertibility and Bijectivity}
\label{subsec:invertibility-bijectivity}
%%%%%%%%%%%%%%%%%%%%%%%%%%%%%%%%%%%%%%%%%%%%%%%%%%%%%%%%%%%%

A function has a genuine inverse function exactly when it is bijective.

If a function is not injective, two different inputs produce the same output,
so the inverse would not know which input to return.

If a function is not surjective, some elements of the codomain have no
preimage, so the inverse could not assign them outputs.

Thus bijectivity is precisely the condition required for invertibility.

%%%%%%%%%%%%%%%%%%%%%%%%%%%%%%%%%%%%%%%%%%%%%%%%%%%%%%%%%%%%
\subsection{Restricting a Domain to Obtain an Inverse}
\label{subsec:domain-restriction-inverse}
%%%%%%%%%%%%%%%%%%%%%%%%%%%%%%%%%%%%%%%%%%%%%%%%%%%%%%%%%%%%

The function

\[ f:\R\to[0,\infty),\qquad f(x)=x^2 \]

is surjective but not injective.

However, restricting its domain gives

\[ f:[0,\infty)\to[0,\infty),\qquad f(x)=x^2. \]

This restricted function is bijective.

Its inverse is

\[ f^{-1}(x)=\sqrt x. \]

%%%%%%%%%%%%%%%%%%%%%%%%%%%%%%%%%%%%%%%%%%%%%%%%%%%%%%%%%%%%
\subsection{Images of Sets}
\label{subsec:images-of-sets}
%%%%%%%%%%%%%%%%%%%%%%%%%%%%%%%%%%%%%%%%%%%%%%%%%%%%%%%%%%%%

Functions can act on entire subsets of their domains.

If

\[ S\subseteq A \]

and

\[ f:A\to B, \]

then the image of \(S\) is

\[ f(S)=\{f(x):x\in S\}. \]

This is read

\begin{center}
``the image of \(S\) under \(f\).''
\end{center}

%%%%%%%%%%%%%%%%%%%%%%%%%%%%%%%%%%%%%%%%%%%%%%%%%%%%%%%%%%%%
\subsection{Worked Example: Image of a Set}
\label{subsec:worked-image-set}
%%%%%%%%%%%%%%%%%%%%%%%%%%%%%%%%%%%%%%%%%%%%%%%%%%%%%%%%%%%%

\begin{workedexamplebox}{Find the Image of a Set}

Let

\[ f(x)=x^2 \]

and

\[ S=\{-2,-1,0,1,2\}. \]

\begin{tabularx}{\textwidth}{@{}>{\raggedright\arraybackslash}p{0.42\textwidth}X@{}}
Evaluate the elements. & \(4,1,0,1,4\)\\
Remove repeated outputs because the image is a set. & \(\{0,1,4\}\)
\end{tabularx}

\begin{answerbox}
\[ f(S)=\{0,1,4\} \]
\end{answerbox}
\end{workedexamplebox}

%%%%%%%%%%%%%%%%%%%%%%%%%%%%%%%%%%%%%%%%%%%%%%%%%%%%%%%%%%%%
\subsection{Preimages of Sets}
\label{subsec:preimages-of-sets}
%%%%%%%%%%%%%%%%%%%%%%%%%%%%%%%%%%%%%%%%%%%%%%%%%%%%%%%%%%%%

If

\[ T\subseteq B, \]

the preimage of \(T\) under \(f:A\to B\) is

\[ f^{-1}(T)=\{x\in A:f(x)\in T\}. \]

Here the notation \(f^{-1}(T)\) does not require \(f\) to have an inverse
function.

It means the set of all inputs whose outputs lie in \(T\).

%%%%%%%%%%%%%%%%%%%%%%%%%%%%%%%%%%%%%%%%%%%%%%%%%%%%%%%%%%%%
\subsection{Worked Example: Preimage of a Set}
\label{subsec:worked-preimage-set}
%%%%%%%%%%%%%%%%%%%%%%%%%%%%%%%%%%%%%%%%%%%%%%%%%%%%%%%%%%%%

\begin{workedexamplebox}{Find a Preimage}

Let

\[ f:\R\to\R,\qquad f(x)=x^2 \]

and

\[ T=[1,4]. \]

\begin{tabularx}{\textwidth}{@{}>{\raggedright\arraybackslash}p{0.42\textwidth}X@{}}
Require the output to lie in \(T\). & \(1\leq x^2\leq4\)\\
Interpret the lower condition. & \(\abs x\geq1\)\\
Interpret the upper condition. & \(\abs x\leq2\)\\
Combine the conditions. & \(1\leq\abs x\leq2\)
\end{tabularx}

\begin{answerbox}
\[ f^{-1}([1,4])=[-2,-1]\cup[1,2] \]
\end{answerbox}
\end{workedexamplebox}

%%%%%%%%%%%%%%%%%%%%%%%%%%%%%%%%%%%%%%%%%%%%%%%%%%%%%%%%%%%%
\subsection{Image Laws}
\label{subsec:image-laws}
%%%%%%%%%%%%%%%%%%%%%%%%%%%%%%%%%%%%%%%%%%%%%%%%%%%%%%%%%%%%

For subsets

\[ S,T\subseteq A, \]

we always have

\[ f(S\cup T)=f(S)\cup f(T). \]

For intersections,

\[ f(S\cap T)\subseteq f(S)\cap f(T). \]

Equality need not hold unless additional conditions, such as injectivity, are
present.

%%%%%%%%%%%%%%%%%%%%%%%%%%%%%%%%%%%%%%%%%%%%%%%%%%%%%%%%%%%%
\subsection{Preimage Laws}
\label{subsec:preimage-laws}
%%%%%%%%%%%%%%%%%%%%%%%%%%%%%%%%%%%%%%%%%%%%%%%%%%%%%%%%%%%%

Preimages interact especially well with set operations.

For subsets

\[ U,V\subseteq B, \]

we have

\[ f^{-1}(U\cup V)=f^{-1}(U)\cup f^{-1}(V), \]

\[ f^{-1}(U\cap V)=f^{-1}(U)\cap f^{-1}(V), \]

and

\[ f^{-1}(U^c)=\left(f^{-1}(U)\right)^c. \]

These identities become important throughout analysis and topology.

%%%%%%%%%%%%%%%%%%%%%%%%%%%%%%%%%%%%%%%%%%%%%%%%%%%%%%%%%%%%
\subsection{Sequences as Functions}
\label{subsec:sequences-as-functions}
%%%%%%%%%%%%%%%%%%%%%%%%%%%%%%%%%%%%%%%%%%%%%%%%%%%%%%%%%%%%

A sequence can be viewed as a function whose domain is a set of integers,
usually \(\N\).

A sequence may be written

\[ a:\N\to\R. \]

Instead of \(a(n)\), sequence notation usually writes

\[ a_n. \]

Thus,

\[ a_n=a(n). \]

For example,

\[ a_n=\frac1n \]

defines the sequence

\[ 1,\frac12,\frac13,\frac14,\ldots. \]

%%%%%%%%%%%%%%%%%%%%%%%%%%%%%%%%%%%%%%%%%%%%%%%%%%%%%%%%%%%%
\subsection{Multivariable Functions}
\label{subsec:multivariable-functions}
%%%%%%%%%%%%%%%%%%%%%%%%%%%%%%%%%%%%%%%%%%%%%%%%%%%%%%%%%%%%

Functions need not take only one input coordinate.

For example,

\[ f:\R^2\to\R \]

may be defined by

\[ f(x,y)=x^2+y^2. \]

This is read

\begin{center}
``\(f\) maps R squared to R''
\end{center}

and

\begin{center}
``\(f\) of \(x,y\) equals \(x^2+y^2\).''
\end{center}

The input is an ordered pair

\[ (x,y)\in\R^2. \]

%%%%%%%%%%%%%%%%%%%%%%%%%%%%%%%%%%%%%%%%%%%%%%%%%%%%%%%%%%%%
\subsection{Vector-Valued Functions}
\label{subsec:vector-valued-functions}
%%%%%%%%%%%%%%%%%%%%%%%%%%%%%%%%%%%%%%%%%%%%%%%%%%%%%%%%%%%%

A function may also produce several output coordinates.

For example,

\[ F:\R\to\R^2,\qquad F(t)=(\cos t,\sin t). \]

This function assigns each real number \(t\) a point in the plane.

Such functions are studied extensively in vector calculus, differential
equations, and geometry.

%%%%%%%%%%%%%%%%%%%%%%%%%%%%%%%%%%%%%%%%%%%%%%%%%%%%%%%%%%%%
\subsection{Functions Between General Sets}
\label{subsec:functions-general-sets}
%%%%%%%%%%%%%%%%%%%%%%%%%%%%%%%%%%%%%%%%%%%%%%%%%%%%%%%%%%%%

Functions need not involve numbers at all.

For example, suppose \(P\) is a set of people and \(D\) is a set of calendar
dates.

A function

\[ b:P\to D \]

might assign each person their birth date.

Likewise, a function could map graphs to numbers, matrices to matrices, sets
to sets, or functions to other functions.

%%%%%%%%%%%%%%%%%%%%%%%%%%%%%%%%%%%%%%%%%%%%%%%%%%%%%%%%%%%%
\subsection{Common Function Errors}
\label{subsec:common-function-errors}
%%%%%%%%%%%%%%%%%%%%%%%%%%%%%%%%%%%%%%%%%%%%%%%%%%%%%%%%%%%%

\subsubsection{Confusing \(f(x)\) with Multiplication}

The expression

\[ f(x) \]

means the value of the function at \(x\).

It does not generally mean \(f\cdot x\).

%%%%%%%%%%%%%%%%%%%%%%%%%%%%%%%%%%%%%%%%%%%%%%%%%%%%%%%%%%%%
\subsubsection{Ignoring the Domain}
\label{subsec:function-error-domain}
%%%%%%%%%%%%%%%%%%%%%%%%%%%%%%%%%%%%%%%%%%%%%%%%%%%%%%%%%%%%

A formula alone does not completely determine a function.

For example,

\[ \frac1{x-3} \]

is undefined at \(x=3\).

Any function using this formula must exclude \(3\) from its domain.

%%%%%%%%%%%%%%%%%%%%%%%%%%%%%%%%%%%%%%%%%%%%%%%%%%%%%%%%%%%%
\subsubsection{Confusing Range and Codomain}
\label{subsec:function-error-range-codomain}
%%%%%%%%%%%%%%%%%%%%%%%%%%%%%%%%%%%%%%%%%%%%%%%%%%%%%%%%%%%%

The codomain is specified as part of the function.

The range consists only of outputs actually obtained.

They need not be equal.

%%%%%%%%%%%%%%%%%%%%%%%%%%%%%%%%%%%%%%%%%%%%%%%%%%%%%%%%%%%%
\subsubsection{Assuming Every Function Has an Inverse}
\label{subsec:function-error-inverse}
%%%%%%%%%%%%%%%%%%%%%%%%%%%%%%%%%%%%%%%%%%%%%%%%%%%%%%%%%%%%

A function must be bijective to possess an inverse function between its full
domain and codomain.

%%%%%%%%%%%%%%%%%%%%%%%%%%%%%%%%%%%%%%%%%%%%%%%%%%%%%%%%%%%%
\subsubsection{Confusing Inverse Function and Reciprocal}
\label{subsec:function-error-inverse-reciprocal}
%%%%%%%%%%%%%%%%%%%%%%%%%%%%%%%%%%%%%%%%%%%%%%%%%%%%%%%%%%%%

The expressions

\[ f^{-1}(x) \]

and

\[ \frac1{f(x)} \]

usually mean different things.

%%%%%%%%%%%%%%%%%%%%%%%%%%%%%%%%%%%%%%%%%%%%%%%%%%%%%%%%%%%%
\subsubsection{Reversing Composition Order}
\label{subsec:function-error-composition}
%%%%%%%%%%%%%%%%%%%%%%%%%%%%%%%%%%%%%%%%%%%%%%%%%%%%%%%%%%%%

The expression

\[ (g\circ f)(x) \]

means

\[ g(f(x)). \]

Therefore \(f\) acts first.

%%%%%%%%%%%%%%%%%%%%%%%%%%%%%%%%%%%%%%%%%%%%%%%%%%%%%%%%%%%%
\subsection{Worked Examples}
\label{subsec:functions-worked-examples}
%%%%%%%%%%%%%%%%%%%%%%%%%%%%%%%%%%%%%%%%%%%%%%%%%%%%%%%%%%%%

\begin{workedexamplebox}{Determine Whether a Relation Is a Function}

Let

\[ R=\{(1,a),(2,b),(3,b),(4,c)\}. \]

\begin{tabularx}{\textwidth}{@{}>{\raggedright\arraybackslash}p{0.42\textwidth}X@{}}
Check input \(1\). & One output: \(a\)\\
Check input \(2\). & One output: \(b\)\\
Check input \(3\). & One output: \(b\)\\
Check input \(4\). & One output: \(c\)
\end{tabularx}

\begin{answerbox}
Yes. Every input has exactly one output, so \(R\) defines a function.
\end{answerbox}
\end{workedexamplebox}

\begin{workedexamplebox}{Find the Natural Domain}

Find the real natural domain of

\[ f(x)=\frac{\sqrt{x-2}}{x-5}. \]

\begin{tabularx}{\textwidth}{@{}>{\raggedright\arraybackslash}p{0.42\textwidth}X@{}}
Require the radicand to be nonnegative. & \(x-2\geq0\Rightarrow x\geq2\)\\
Require the denominator to be nonzero. & \(x\neq5\)\\
Combine the restrictions. & \([2,5)\cup(5,\infty)\)
\end{tabularx}

\begin{answerbox}
\[ \dom(f)=[2,5)\cup(5,\infty) \]
\end{answerbox}
\end{workedexamplebox}

\begin{workedexamplebox}{Test Injectivity}

Consider

\[ f:\R\to\R,\qquad f(x)=x^3+1. \]

\begin{tabularx}{\textwidth}{@{}>{\raggedright\arraybackslash}p{0.42\textwidth}X@{}}
Assume equal outputs. & \(x_1^3+1=x_2^3+1\)\\
Subtract \(1\). & \(x_1^3=x_2^3\)\\
Take cube roots. & \(x_1=x_2\)
\end{tabularx}

\begin{answerbox}
\(f\) is injective.
\end{answerbox}
\end{workedexamplebox}

%%%%%%%%%%%%%%%%%%%%%%%%%%%%%%%%%%%%%%%%%%%%%%%%%%%%%%%%%%%%
\subsection{Exercises}
\label{subsec:functions-exercises}
%%%%%%%%%%%%%%%%%%%%%%%%%%%%%%%%%%%%%%%%%%%%%%%%%%%%%%%%%%%%

\begin{exercise}[1]
Let
\[ f(x)=3x-4. \]
Find
\[ f(2),\qquad f(-1),\qquad f(0). \]
\end{exercise}

\begin{exercise}[1]
Determine whether
\[ \{(1,a),(2,b),(3,a)\} \]
defines a function.
\end{exercise}

\begin{exercise}[1]
Determine whether
\[ \{(1,a),(1,b),(2,c)\} \]
defines a function.
\end{exercise}

\begin{exercise}[1]
Find the natural real domain of
\[ f(x)=\frac1{x+4}. \]
\end{exercise}

\begin{exercise}[1]
Find the natural real domain of
\[ f(x)=\sqrt{5-x}. \]
\end{exercise}

\begin{exercise}[1]
Let
\[ f(x)=
\begin{cases}
2x+1,&x\geq0,\\
x^2,&x<0.
\end{cases}
\]
Find \(f(3)\), \(f(-2)\), and \(f(0)\).
\end{exercise}

\begin{exercise}[2]
Determine whether
\[ f:\R\to\R,\qquad f(x)=5x+2 \]
is injective.
\end{exercise}

\begin{exercise}[2]
Determine whether
\[ f:\R\to\R,\qquad f(x)=x^2 \]
is surjective.
\end{exercise}

\begin{exercise}[2]
Determine whether
\[ f:\R\to[0,\infty),\qquad f(x)=x^2 \]
is surjective.
\end{exercise}

\begin{exercise}[2]
Determine whether
\[ f:\R\to\R,\qquad f(x)=x^3 \]
is bijective.
\end{exercise}

\begin{exercise}[1]
Let
\[ f(x)=x+2,\qquad g(x)=3x. \]
Find
\[ (g\circ f)(x) \]
and
\[ (f\circ g)(x). \]
\end{exercise}

\begin{exercise}[2]
Find the inverse of
\[ f(x)=5x-9. \]
\end{exercise}

\begin{exercise}[2]
Verify your inverse from the previous exercise by computing
\[ f^{-1}(f(x)) \]
and
\[ f(f^{-1}(x)). \]
\end{exercise}

\begin{exercise}[2]
Let
\[ f(x)=x^2 \]
and
\[ S=\{-3,-2,-1,0,1,2,3\}. \]
Find \(f(S)\).
\end{exercise}

\begin{exercise}[2]
For
\[ f(x)=x^2, \]
find
\[ f^{-1}([4,9]). \]
\end{exercise}

\begin{exercise}[2]
Give an example of a function that is injective but not surjective.
Clearly specify its domain and codomain.
\end{exercise}

\begin{exercise}[2]
Give an example of a function that is surjective but not injective.
Clearly specify its domain and codomain.
\end{exercise}

\begin{exercise}[2]
Explain why the codomain matters when deciding whether a function is
surjective.
\end{exercise}

\begin{exercise}[3]
Prove that the composition of two injective functions is injective.
\end{exercise}

\begin{exercise}[3]
Prove that the composition of two surjective functions is surjective.
\end{exercise}

\begin{exercise}[3]
Prove that the composition of two bijective functions is bijective.
\end{exercise}

\begin{exercise}[3]
Prove
\[ f^{-1}(U\cup V)=f^{-1}(U)\cup f^{-1}(V). \]
\end{exercise}

\begin{exercise}[3]
Prove
\[ f^{-1}(U\cap V)=f^{-1}(U)\cap f^{-1}(V). \]
\end{exercise}

%%%%%%%%%%%%%%%%%%%%%%%%%%%%%%%%%%%%%%%%%%%%%%%%%%%%%%%%%%%%
\subsection{Section Connections}
\label{subsec:functions-connections}
%%%%%%%%%%%%%%%%%%%%%%%%%%%%%%%%%%%%%%%%%%%%%%%%%%%%%%%%%%%%

\chapterconnection{
Functions build directly on relations from \cref{sec:relations} and become
foundational throughout the rest of mathematics. Calculus studies rates of
change and accumulation of functions; linear algebra studies linear
transformations; probability studies random variables; statistics studies
models and estimators; topology and analysis rely heavily on images and
preimages; and abstract algebra studies functions preserving algebraic
structure.
}

%%%%%%%%%%%%%%%%%%%%%%%%%%%%%%%%%%%%%%%%%%%%%%%%%%%%%%%%%%%%
% SECTION SOURCE TRACKING
%%%%%%%%%%%%%%%%%%%%%%%%%%%%%%%%%%%%%%%%%%%%%%%%%%%%%%%%%%%%

%------------------------------------------------------------
% 0.6 Functions
%------------------------------------------------------------
%
% Sources:
%   - Richard Hammack, Book of Proof, Third Edition.
%     Key: hammack2018bookofproof
%
% Used for:
%   - functions as special relations;
%   - domain, codomain, and image terminology;
%   - injective, surjective, and bijective functions;
%   - function composition;
%   - inverse functions;
%   - images and preimages;
%   - standard proof conventions involving functions.
%
% Notes:
%   - Algebraic and calculus-specific function families are developed later.
%   - Graphing techniques are deferred to Analysis 5.1.
%   - Linear transformations receive full treatment in Linear Algebra.
%
\nocite{hammack2018bookofproof}